\section{Pin Connections}
\label{app:pin_connections}

\begin{table}[H]
    \centering
    \caption{External pin summary.}
    \label{tab:appendix_pins}
    \renewcommand{\arraystretch}{1.2}
    \small
    \begin{tabularx}{0.96\textwidth}{p{3.4cm} p{3.5cm} X}
        \toprule
        \textbf{Connection} & \textbf{DE2 GPIO} & \textbf{Description} \\
        \midrule
        Sensor 1 SCL/SDA & GPIO[0]/GPIO[1] & Independent I2C bus for magnetometer 1. \\
        Sensor 2 SCL/SDA & GPIO[2]/GPIO[3] & Independent I2C bus for magnetometer 2. \\
        Sensor 3 SCL/SDA & GPIO[4]/GPIO[5] & Independent I2C bus for magnetometer 3. \\
        DRV8825 STEP & GPIO[8] & Step pulse output. \\
        DRV8825 DIR & GPIO[9] & Direction output. \\
        DRV8825 ENABLE\_N & GPIO[10] & Active-low driver enable. \\
        75 Hz press button & GPIO[11] & Active-high press state for 75 Hz finger. \\
        45 Hz press button & GPIO[12] & Active-high press state for 45 Hz finger. \\
        \bottomrule
    \end{tabularx}
\end{table}

\section{UART Frame Reference}
\label{app:uart_reference}

\begin{lstlisting}
H2,75,H75S1,H75S2,H75S3,45,H45S1,H45S2,H45S3
KEY,75,KEY75,PRESS75,45,KEY45,PRESS45
\end{lstlisting}

The $H^2$ values are unsigned 32-bit hexadecimal Q16 values. The key indices are decimal global key indices. The press bits are active-high digital states after FPGA debouncing.

\section{Important Runtime Parameters}
\label{app:runtime_parameters}

\begin{table}[H]
    \centering
    \caption{Important final parameters.}
    \label{tab:runtime_parameters}
    \renewcommand{\arraystretch}{1.2}
    \begin{tabular}{l c}
        \toprule
        \textbf{Parameter} & \textbf{Value} \\
        \midrule
        Magnetometer count & 3 \\
        Magnetometer output data rate & 200 Hz \\
        Magnetometer range & $\pm 2$ G \\
        Finger frequencies & 75 Hz and 45 Hz \\
        Lock-in window & 100 samples \\
        Classifier $H^2$ average & 10 frames \\
        LUT entries per frequency & 45 \\
        Local key range & -2 to +2 \\
        Global key range & 0 to 14 \\
        One-key movement & 100 motor steps \\
        Stepper speed & 800 steps/s \\
        Step pulse width & 5 us \\
        Platform cooldown & 100 ms \\
        \bottomrule
    \end{tabular}
\end{table}

\section{Code Organization}
\label{app:code_organization}

The main FPGA source files are located in \code{src/}. The main Python utility files are located in \code{src\_python/}. The final PC piano interface is \code{src\_python/piano\_key\_monitor.py}. The FPGA LUT ROM files are \code{src/h75\_lut\_3sensor.mem} and \code{src/h45\_lut\_3sensor.mem}.

\newpage
\section{Project File Structure}
\label{app:file_structure}

The submitted repository is organized as follows:

\begin{lstlisting}
final_project/
|-- README.md
|-- src/
`-- src_python/
\end{lstlisting}

The Verilog/SystemVerilog files in \code{src/} implement the FPGA runtime system. The Python files in \code{src\_python/} are used for LUT collection, LUT export, debugging, and the final PC piano interface. The final report source and generated PDF are stored in \code{report/latex/}.
